\documentclass[11pt]{article}
\usepackage{amsmath,amssymb,geometry}
\geometry{margin=24mm}
\newcommand{\E}{\mathbb E}
\newcommand{\dd}{\,\mathrm d}
\title{G371: a rotation control with the shifted soft propagator retained}
\author{Phi3Lab development record}
\date{25 October 2026}
\begin{document}
\maketitle
\noindent This is an exact change of the integration density. It is not a new counterterm or a precision certificate. All four forests of the original graph, its local Taylor orders, and its overall zeroth-order subtraction remain unchanged. The mass is normalized to one and the external Gram matrix is
\[
 p_1^2=p_2^2=1,\qquad p_1\cdot p_2=-\tfrac12,\qquad P=p_1+p_2,\quad P^2=1.
\]
The original-program SHA256 is
\begin{center}\ttfamily\small
55c06ff5efb4fc30320d30378e00f9e94482\\
b2ad57de4a9023f14e192ea7730a.
\end{center}

\section*{An exact denominator in the control}
Set \(S=q+r\), \(t=q^2\), \(v=r^2\), \(h=q\cdot r\),
\(w=S^2=t+v+2h\), \(X=1+w\), and \(\ell=2S\cdot P\).
Separate the two cross propagators from the other four host propagators and the two complete triangle factors. Relative to the four unshifted host denominators, their product has the physical external-momentum expansion
\[
 F(\lambda)=F_0+\lambda F_1+\lambda^2F_2+\cdots,\qquad
 Y(\lambda)=X+\lambda\ell+\lambda^2.
\]
The sign average of the complete subtracted numerator is
\[
 D=\frac12\sum_{\sigma=\pm1}\frac{F(\sigma)}{X Y(\sigma)}
       -\frac{F_0}{X^2}.
\]
Use the control
\[
 A=\frac12\sum_{\sigma=\pm1}
 \frac{F_0+\sigma F_1+F_2}{X Y(\sigma)}-\frac{F_0}{X^2}.
\]
Then replace \(D\) by \(D-A+\E_{\rm rot} A\), at fixed \(t,v,h\).
The cross denominator \(Y(\sigma)\) is never Taylor expanded. This matters when
\(q,r\) are large and \(S\) is small. The remaining point residual is
\[
 D-A=\frac12\sum_{\sigma=\pm1}
 \frac{F(\sigma)-F_0-\sigma F_1-F_2}{X Y(\sigma)}.
\]
The four baseline host denominators contribute the common factor
\(D_4=(1+t)^2(1+v)^2\).

\section*{Weighted rotation moments}
For \(\mu=\widehat S\cdot P\), the normalized six-dimensional measure is
\[
 \dd\nu(\mu)=\frac8{3\pi}(1-\mu^2)^{3/2}\dd\mu,\qquad -1<\mu<1.
\]
Put \(a=w+2\), \(b=2\sqrt w\), \(R=\sqrt{w^2+4}\), and \(s=a+R\).
The needed moments are
\begin{align*}
 M&=\E\frac1{a+b\mu}=\frac{4(a+2R)}{3s^2},\\
 M_1&=\E\frac{\mu}{a+b\mu}=-\frac{b(a+3R)}{3s^3},\\
 M_2&=\E\frac{\mu^2}{a+b\mu}=\frac{a(a+3R)}{3s^3}.
\end{align*}
They satisfy \(aM+bM_1=1\), \(aM_1+bM_2=0\), and have limits
\((M,M_1,M_2)=(1/2,0,1/12)\) at \(w=0\).

Decompose
\[
 q=\alpha S+k,\quad r=\beta S-k,\quad
 \alpha=\frac{t+h}{w},\quad \beta=\frac{v+h}{w},\quad
 k\cdot S=0,\quad k^2=\rho^2=\frac{tv-h^2}{w}.
\]
At fixed \(S\), \(k\) is uniform in its five-dimensional perpendicular space.
Write \(q_0=q,q_1=r\), \(c_0=t+h,c_1=v+h\), and \(\epsilon_0=1,\epsilon_1=-1\).
For \(G_{ab}=p_a\cdot p_b\), define
\[
 V_{ab}=\frac{M_2}{4}+
       \left(G_{ab}-\frac14\right)\frac{M-M_2}{5}.
\]
The first and second weighted projection moments are
\begin{align*}
 \E\frac{q_i\cdot p_a}{Y(1)}
   &=-\frac{c_i(a+3R)}{3s^3},\\
 \E\frac{(q_i\cdot p_a)(q_j\cdot p_b)}{Y(1)}
   &=\frac{c_i c_j}{w}V_{ab}
    +\frac{\epsilon_i\epsilon_j(tv-h^2)}{5w}
       \left(G_{ab}M-V_{ab}\right).
\end{align*}
If \(\overline F_1,\overline F_2\) denote the corresponding weighted linear
and quadratic averages, then
\[
 \E_{\rm rot}A
 =F_0\frac{XM-1}{X^2}
  +\frac{\overline F_1+\overline F_2}{X}.
\]
The \(w=0\) boundary is treated by limits; sampled coordinates lie in the
strict Gram domain. No Gram clipping is used.

\section*{Intrinsic angles and normalized return density}
An explicit frame is
\begin{align*}
 e&=(\sqrt z\cos\theta,\sqrt z\sin\theta,\sqrt{1-z},0),\\
 n_1&=(-\sin\theta,\cos\theta,0,0),\\
 n_2&=(\sqrt{1-z}\cos\theta,\sqrt{1-z}\sin\theta,-\sqrt z,0),\\
 n_3&=(0,0,0,1),\\
 n&=\sqrt Z(\cos\phi\,n_1+\sin\phi\,n_2)+\sqrt{1-Z}\,n_3,\\
 q&=\sqrt t\,e,\qquad r=\sqrt v(\cos\psi\,e+\sin\psi\,n).
\end{align*}
Here \(\theta,\phi\) each cover \(2\pi\), and \(\psi=\pi u\).
The overall-rotation factor in unit-cube coordinates is
\(O=3(1-z)\sqrt{1-Z}\), whose integral over \(z,Z,\theta/(2\pi),\phi/(2\pi)\)
is one. The intrinsic-angle factor is \(I_\psi=(8/3)\sin^4(\pi u)\).
Thus the returned density is
\[
 {\cal J}_{\rm rad} I_\psi
 \left[ O\,\frac{D-A}{D_4}+\frac{\E_{\rm rot}A}{D_4}\right].
\]
For \(t=[s_0/(1-s_0)]^2,v=[s_1/(1-s_1)]^2\),
\[
 {\cal J}_{\rm rad}=\frac{t^3v^3}{s_0(1-s_0)s_1(1-s_1)}.
\]
For \(T=[s_0/(1-s_0)]^2,t=T s_1,v=T(1-s_1)\),
\[
 {\cal J}_{\rm rad}=\frac{T^6s_1^2(1-s_1)^2}{2s_0(1-s_0)}.
\]
Both contain the original \(t^2v^2\dd t\dd v/4\) measure.
The three unimodular charts \((u,v),(u,v-u),(v-u,u)\) retain their
rotation-invariant weights \(D_{\rm excluded}^3/
(D_q^3+D_r^3+D_{q+r}^3)\). Their sum is one; no extra factor \(1/3\) is inserted.

\section*{Absolute integrability and independent checking}
With \(M_0=256\log(2+\max(t,v,w))\) and
\(D_{\min}=1+\min(t,v,w)\), the frozen route proof gives
\[
 |F_0|\le M_0^2,\quad
 |F_1|\le 81\cdot18\,M_0^2D_{\min}^{-1/2},\quad
 |F_2|\le81\cdot206\,M_0^2D_{\min}^{-1}.
\]
The even difference, odd part, and even part of \(X/Y\) are bounded by
\(57/X,18/\sqrt X,3\), respectively. Consequently the control's relative
numerator is bounded by \(76359M_0^2/D_{\min}\). The complete residual plus
returned mean is bounded in absolute integral by
\[
 (729\cdot808/2+2\cdot76359)\frac{M_0^2}{D_{\min}}
 =447234\frac{M_0^2}{D_{\min}},
\]
below the previously proved integrable bound
\(883548M_0^2/D_{\min}\), with the same baseline host denominators.

The native implementation uses the explicit covariance above. An independent
reference extracts the degree-two polynomial from formal Taylor convolution,
uses ten directions for \(S\) at fixed \(\mu\) and ten perpendicular directions
for \(k\), then reconstructs its quadratic conditional mean in \(\mu\).
Direct angular and child-parameter integrations supply a further check.
Proofs, pointwise checks, fixed diagnostics, independent production, and replay
audits are separate gates. Their actual exits and hashes are recorded in the
checkpoint; this derivation asserts no new accepted integral.
\end{document}